\chapter{A Concrete Naming Certificate for $\GroundCompilerRecognizerAuditPacketUp$}
\label{ch:concrete-instances-groundcompilerrecognizerauditpacket-namecert}
\origin{human}

The $\GroundCompilerRecognizerAuditPacketUp$ carrier turns the recognizer-facing part of the GroundCompiler audit map of \autoref{ch:bedc-groundcompiler-audit-map} into a finite BEDC packet. It sits beside the audit-map family surface of \autoref{ch:concrete-instances-auditmapfamily-namecert}, the certificate-compiler graph packet of \autoref{ch:concrete-instances-certificatecompiler-namecert}, and the taste-gate compiler witness of \autoref{ch:concrete-instances-tastegatecompilerwitness-namecert}. Those sibling chapters supply the family template, certificate handoff route, and chapter-gate witness rows; this chapter packages the GroundCompiler recognizer rows and non-evidence rows as the object those consumers may inspect.

\begin{definition}[GroundCompiler recognizer audit packet carrier]
\label{def:groundcompiler-recognizer-audit-packet-carrier}
A $\GroundCompilerRecognizerAuditPacketUp$ carrier is a finite $\BHist$ packet
$$
G=(R,A,E,Q,B,H,C,P,N)
$$
with the following displayed rows:
\begin{enumerate}
\item $R$ is the recognizer row naming the generated recognizers, accepted exports, and recognizer-generatedness checks from the GroundCompiler audit map;
\item $A$ is the address-analysis row recording compiler-layer address, theorem-code, chapter-code, and classifier-quotient reads that are permitted only as displayed audit data;
\item $E$ is the evidence row containing event-flow conservativity, no-hidden-input, channel-bijection, and source-channel separation rows available to certificate consumers;
\item $Q$ is the non-evidence row listing code-not-proof, report-not-evidence, motif-not-identity, quotient-not-raw-equality, and cannot-claim boundaries;
\item $B$ is the bridge-request row that records which downstream packet may consume the recognizer audit surface, including the sibling packets of \autoref{ch:concrete-instances-auditmapfamily-namecert}, \autoref{ch:concrete-instances-certificatecompiler-namecert}, and \autoref{ch:concrete-instances-tastegatecompilerwitness-namecert};
\item $H$ records componentwise $\hsame$ transports for the recognizer, address-analysis, evidence, non-evidence, bridge-request, route, provenance, and local naming rows;
\item $C$ records the displayed $\Cont$ routes by which a consumer inspects $R,A,E,Q,B$ without adding a hidden compiler proof row;
\item $P$ is the $\Pkg$ provenance over $\GroundCompiler$, \autoref{ch:bedc-groundcompiler-audit-map}, the sibling concrete packets, $\BHist$, $\hsame$, $\Cont$, $\Pkg$, and $\NameCert$;
\item $N$ is the local $\NameCert_{\GroundCompilerRecognizerAuditPacketUp}$ row.
\end{enumerate}
The carrier accepts exactly the displayed recognizer, address-analysis, evidence, non-evidence, bridge-request, transport, route, provenance, and local naming rows. It exports no proof that code is proof, no report-to-evidence conversion, no raw classifier equality, no hidden self-hosting tower, and no MetaCIC theorem outside the listed sibling handoff.
\end{definition}
\leanchecked{BEDC.Derived.GroundCompilerRecognizerAuditPacketUp.GroundCompilerRecognizerAuditPacketTasteGate\_single\_carrier\_alignment}

\begin{theorem}[GroundCompiler recognizer audit packet obligation surface]
\label{thm:groundcompiler-recognizer-audit-packet-obligation-surface}
For the carrier of \autoref{def:groundcompiler-recognizer-audit-packet-carrier}, the local $\NameCert_{\GroundCompilerRecognizerAuditPacketUp}$ surface consists exactly of carrier admission for
$$
G=(R,A,E,Q,B,H,C,P,N),
$$
recognizer readback through $R$, address-analysis readback through $A$, evidence readback through $E$, non-evidence readback through $Q$, bridge-request readback through $B$, componentwise transport through $H$, displayed continuation through $C$, package provenance through $P$, and local naming through $N$.
\end{theorem}
\begin{proof}
Carrier admission is tuple projection from \autoref{def:groundcompiler-recognizer-audit-packet-carrier}. The recognizer, address-analysis, evidence, non-evidence, and bridge-request rows are all displayed before any downstream consumer can inspect the packet.

The transport and route rows preserve that displayed inventory. Componentwise $\hsame$ entries in $H$ can only move the named rows, and the continuations in $C$ can only replay routes through $R,A,E,Q,B$.

Finally $P$ and $N$ account for provenance and local naming. Since provenance cites the GroundCompiler audit map and the sibling concrete packets only as row sources, the local NameCert surface is exhausted by the nine displayed coordinates.
\end{proof}

\begin{theorem}[GroundCompiler recognizer audit packet non-evidence boundary]
\label{thm:groundcompiler-recognizer-audit-packet-non-evidence-boundary}
Every public read of an accepted $\GroundCompilerRecognizerAuditPacketUp$ carrier keeps the non-evidence row $Q$ visible. A consumer may cite recognizer-generatedness, accepted exports, event-flow conservativity, no-hidden-input, channel separation, and address-analysis rows only together with the boundaries that code is not proof, reports are not evidence, motifs are not object identity, classifier quotients are not raw equality, and conditional self-hosting is not an unconditional compiler tower.
\end{theorem}
\begin{proof}
Start with the carrier projection. The positive recognizer and compiler-facing rows are $R,A,E$, while the refusal row is $Q$.

Apply the obligation surface of \autoref{thm:groundcompiler-recognizer-audit-packet-obligation-surface}. Any route that exposes $R,A,E$ also exposes the same packet through $H,C,P,N$, so the route cannot drop $Q$ without leaving the carrier.

The row $Q$ is therefore part of the public read, not a side note. Since the packet contains no coordinate converting code, reports, motifs, quotients, or conditional self-hosting into stronger evidence, every consumer remains inside the listed non-evidence boundary.
\end{proof}

\begin{theorem}[GroundCompiler recognizer audit packet consumer handoff]
\label{thm:groundcompiler-recognizer-audit-packet-consumer-handoff}
The bridge-request row $B$ hands an accepted $\GroundCompilerRecognizerAuditPacketUp$ packet only to displayed BEDC consumers. In particular, audit-map family consumers read the packet as a family row, certificate-compiler consumers read it as finite compiler-side provenance, and taste-gate compiler witnesses read it as chapter-gate evidence tied to the same recognizer and non-evidence rows.
\end{theorem}
\begin{proof}
Project the bridge-request coordinate $B$. Its listed sibling destinations are the audit-map family, certificate compiler, and taste-gate compiler witness chapters named in the carrier provenance.

For the audit-map family, the handoff lands as a visible family row and keeps the positive, conditional, obstruction, and frontier split supplied by \autoref{ch:concrete-instances-auditmapfamily-namecert}. For the certificate compiler, the handoff supplies finite compiler-side provenance rather than host compiler semantics. For the taste-gate compiler witness, the handoff supplies named chapter-gate evidence rather than an ambient approval principle.

The route rows $H$ and $C$ preserve the same $R,A,E,Q,B$ inventory during each handoff. Thus no downstream consumer receives a stronger object than the displayed recognizer audit packet.
\end{proof}

\begin{theorem}[GroundCompiler recognizer audit packet NameCert]
\label{thm:groundcompiler-recognizer-audit-packet-namecert}
The public $\NameCert_{\GroundCompilerRecognizerAuditPacketUp}$ packet is exactly the finite surface
$$
R,A,E,Q,B,H,C,P,N.
$$
It packages the recognizer and non-evidence rows of \autoref{ch:bedc-groundcompiler-audit-map} for concrete BEDC consumers, and it exports no hidden compiler proof, no proof-producing code interpretation, no evidence extracted from reports, no raw equality behind classifier quotients, and no unstated MetaCIC closure.
\end{theorem}
\begin{proof}
Read the public packet by carrier projection and the local obligation surface. The available rows are the recognizer row $R$, address-analysis row $A$, evidence row $E$, non-evidence row $Q$, bridge-request row $B$, transports $H$, routes $C$, provenance $P$, and local naming $N$.

Use \autoref{thm:groundcompiler-recognizer-audit-packet-non-evidence-boundary} to keep the refusal rows attached to every recognizer-facing read. Use \autoref{thm:groundcompiler-recognizer-audit-packet-consumer-handoff} to restrict downstream use to the displayed sibling consumers.

No additional coordinate remains. Therefore the public NameCert packet is exactly $R,A,E,Q,B,H,C,P,N$, with the listed compiler-proof, report-evidence, raw-equality, and MetaCIC-closure claims outside the carrier.
\end{proof}

\begin{closurestatus}{\GroundCompilerRecognizerAuditPacketUp}
  \constructivestory{A $\GroundCompilerRecognizerAuditPacketUp$ packet is generated from finite $\BHist$ rows for recognizers, address analysis, positive evidence, non-evidence boundaries, bridge requests, componentwise transports, displayed continuations, provenance, and local naming. Recognizer-facing reads are replayed by $\Cont$ through $R,A,E,Q,B$, so evidence rows and refusal rows remain attached in one public packet. $\hsame$ transports the recognizer, address, evidence, non-evidence, and bridge rows without strengthening them, while $\Pkg$ records the GroundCompiler audit map and sibling concrete consumers. The packet packages recognizer audit data without turning code or reports into proofs.}
  \theoryclosure{\seedClosure}
  \scopeclosed{This chapter binds the GroundCompiler recognizer audit packet carrier, obligation surface, non-evidence boundary, consumer handoff, and public NameCert packet.}
  \formalstatus{\unformalizedV}
  \bridgestatus{none}
  \origin{human}
  \notclaimed{It does not claim compiler proofs, report-to-evidence conversion, raw classifier equality, hidden self-hosting towers, or unstated MetaCIC closure.}
  \upgradepath{Obligation closure requires Lean encodings of the recognizer audit carrier, non-evidence preservation, consumer handoff, and public packet theorem.}
\end{closurestatus}
